% --- Archon physics formalization source begin ---
% archon:physics
% archon:covers PhyXMiniProblems/problem_phyx_mini_0831.lean
% archon:source-report reports/phyx_mini/problem_phyx_mini_0831.source.json

\chapter{Physics problem phyx\_mini\_0831}
\label{ch:PhyXMiniProblems_problem_phyx_mini_0831}

\paragraph{Problem source.}
\#\# Question

What is the strength of the electric field at the position indicated by the dot?

\#\# Answer choices

- A: A: 1.3 \textbackslash{}times 10\textasciicircum{}\{3\}N/C

- B: B: 5.3 \textbackslash{}times 10\textasciicircum{}\{3\}N/C

- C: C: 1.0 \textbackslash{}times 10\textasciicircum{}\{4\}N/C

- D: D: 7.6 \textbackslash{}times 10\textasciicircum{}\{3\}N/C

\#\# Figure caption (auxiliary; use the image as primary evidence)

The image shows a diagram with three main elements:

1. **Charges**:
   - A red circle with a "+", labeled "3.0 nC" positioned at the top.
   - A blue circle with a "-", labeled "-3.0 nC" positioned below.

2. **Distances**:
   - A vertical distance of "5.0 cm" between the positive charge (red circle) and an unspecified point between two horizontal arrows.
   - Another vertical distance of "5.0 cm" between the unspecified point and the negative charge (blue circle).
   - A horizontal arrow extending from the unspecified central point, labeled "5.0 cm".

3. **Central Point**:
   - Marked with a small black dot where the horizontal arrow meets, serving as a reference point for the distances from the charges.

The setup resembles a configuration for calculating electric field or potential due to point charges.

\#\# Dataset metadata

Electrostatics; Physical Model Grounding Reasoning, Multi-Formula Reasoning

\paragraph{Recorded answer/context.}
D: D: 7.6 \textbackslash{}times 10\textasciicircum{}\{3\}N/C

\paragraph{Figure/image path.}
/root/proposal\_for\_physic/science-mango/phyx\_mini\_run/phyx\_data/test\_image/831.png

\paragraph{Formalization target.}
create a compiling Lean file with sorry bodies at `PhyXMiniProblems/problem\_phyx\_mini\_0831.lean`. The Lean declarations must preserve the physical quantities, dimensions or dimensional roles, figure labels, governing-law hypotheses, and final relation expressed by this problem.
Use Mathlib/Physlib names found through LeanExplore where available. If a physics API is missing, introduce faithful local abstractions rather than scalar placeholder aliases.

\begin{theorem}[Physics formalization target]
\label{thm:physics:phyx_mini_0831:target}
\lean{PhyXMiniProblems.ProblemPhyXMini0831.electric_field_strength_at_observation_dot}
\uses{def:physics:phyx-mini-0831:phyxminiproblems-problemphyxmini0831-oppositepointchargesetup, def:physics:phyx-mini-0831:phyxminiproblems-problemphyxmini0831-matchesprimaryfigure, def:physics:phyx-mini-0831:phyxminiproblems-problemphyxmini0831-matchesproblemreadouts, def:physics:phyx-mini-0831:phyxminiproblems-problemphyxmini0831-hasdepictedgeometry, def:physics:phyx-mini-0831:phyxminiproblems-problemphyxmini0831-hasphysicalparameters, def:physics:phyx-mini-0831:phyxminiproblems-problemphyxmini0831-satisfiespointchargeelectrostatics, def:physics:phyx-mini-0831:phyxminiproblems-problemphyxmini0831-answeragreementtoleranceinnewtonspercoulomb, def:physics:phyx-mini-0831:phyxminiproblems-problemphyxmini0831-answerdiscrepancyinnewtonspercoulomb, def:physics:phyx-mini-0831:phyxminiproblems-problemphyxmini0831-isuniqueclosestanswer}
The assigned autoformalize agent should translate this physics problem into Lean declarations in the covered file, with theorem and lemma proof bodies written as `by sorry`.
\end{theorem}
\begin{proof}
This is an autoformalization task, not a proof task. Produce statements that can later be proved by the physics prover without weakening the physical model.
\end{proof}

% --- Archon declaration topology begin ---
\section{Declaration topology}
\begin{definition}[electric Field Dimension]
  \label{def:physics:phyx-mini-0831:phyxminiproblems-problemphyxmini0831-electricfielddimension}
  \lean{PhyXMiniProblems.ProblemPhyXMini0831.electricFieldDimension}
  The physical dimension M L T⁻² C⁻¹ of electric-field strength
\end{definition}
\begin{proof}
  This is the declared model definition of electric Field Dimension.
\end{proof}
\begin{definition}[Length Quantity]
  \label{def:physics:phyx-mini-0831:phyxminiproblems-problemphyxmini0831-lengthquantity}
  \lean{PhyXMiniProblems.ProblemPhyXMini0831.LengthQuantity}
  A nonnegative unit-independent physical length
\end{definition}
\begin{proof}
  This is the declared model definition of Length Quantity.
\end{proof}
\begin{definition}[Charge Quantity]
  \label{def:physics:phyx-mini-0831:phyxminiproblems-problemphyxmini0831-chargequantity}
  \lean{PhyXMiniProblems.ProblemPhyXMini0831.ChargeQuantity}
  A signed unit-independent electric charge
\end{definition}
\begin{proof}
  This is the declared model definition of Charge Quantity.
\end{proof}
\begin{definition}[Planar Position Quantity]
  \label{def:physics:phyx-mini-0831:phyxminiproblems-problemphyxmini0831-planarpositionquantity}
  \lean{PhyXMiniProblems.ProblemPhyXMini0831.PlanarPositionQuantity}
  A signed position vector in the two-dimensional plane of the figure
\end{definition}
\begin{proof}
  This is the declared model definition of Planar Position Quantity.
\end{proof}
\begin{definition}[Planar Electric Field Quantity]
  \label{def:physics:phyx-mini-0831:phyxminiproblems-problemphyxmini0831-planarelectricfieldquantity}
  \lean{PhyXMiniProblems.ProblemPhyXMini0831.PlanarElectricFieldQuantity}
  \uses{def:physics:phyx-mini-0831:phyxminiproblems-problemphyxmini0831-electricfielddimension}
  A physical electric-field vector in the plane of the figure
\end{definition}
\begin{proof}
  This is the declared model definition of Planar Electric Field Quantity.
\end{proof}
\begin{definition}[Electric Field Strength Quantity]
  \label{def:physics:phyx-mini-0831:phyxminiproblems-problemphyxmini0831-electricfieldstrengthquantity}
  \lean{PhyXMiniProblems.ProblemPhyXMini0831.ElectricFieldStrengthQuantity}
  \uses{def:physics:phyx-mini-0831:phyxminiproblems-problemphyxmini0831-electricfielddimension}
  A nonnegative physical electric-field strength
\end{definition}
\begin{proof}
  This is the declared model definition of Electric Field Strength Quantity.
\end{proof}
\begin{definition}[x Axis]
  \label{def:physics:phyx-mini-0831:phyxminiproblems-problemphyxmini0831-xaxis}
  \lean{PhyXMiniProblems.ProblemPhyXMini0831.xAxis}
  Coordinate 0, horizontal and positive to the right
\end{definition}
\begin{proof}
  This is the declared model definition of x Axis.
\end{proof}
\begin{definition}[y Axis]
  \label{def:physics:phyx-mini-0831:phyxminiproblems-problemphyxmini0831-yaxis}
  \lean{PhyXMiniProblems.ProblemPhyXMini0831.yAxis}
  Coordinate 1, vertical and positive upward
\end{definition}
\begin{proof}
  This is the declared model definition of y Axis.
\end{proof}
\begin{definition}[length In Meters]
  \label{def:physics:phyx-mini-0831:phyxminiproblems-problemphyxmini0831-lengthinmeters}
  \lean{PhyXMiniProblems.ProblemPhyXMini0831.lengthInMeters}
  \uses{def:physics:phyx-mini-0831:phyxminiproblems-problemphyxmini0831-lengthquantity}
  Read a physical length in coherent-SI metres
\end{definition}
\begin{proof}
  This is the declared model definition of length In Meters.
\end{proof}
\begin{definition}[length In Centimeters]
  \label{def:physics:phyx-mini-0831:phyxminiproblems-problemphyxmini0831-lengthincentimeters}
  \lean{PhyXMiniProblems.ProblemPhyXMini0831.lengthInCentimeters}
  \uses{def:physics:phyx-mini-0831:phyxminiproblems-problemphyxmini0831-lengthquantity, def:physics:phyx-mini-0831:phyxminiproblems-problemphyxmini0831-lengthinmeters}
  Read a physical length in centimetres
\end{definition}
\begin{proof}
  This is the declared model definition of length In Centimeters.
\end{proof}
\begin{definition}[charge In Coulombs]
  \label{def:physics:phyx-mini-0831:phyxminiproblems-problemphyxmini0831-chargeincoulombs}
  \lean{PhyXMiniProblems.ProblemPhyXMini0831.chargeInCoulombs}
  \uses{def:physics:phyx-mini-0831:phyxminiproblems-problemphyxmini0831-chargequantity}
  Read a signed physical charge in coherent-SI coulombs
\end{definition}
\begin{proof}
  This is the declared model definition of charge In Coulombs.
\end{proof}
\begin{definition}[charge In Nanocoulombs]
  \label{def:physics:phyx-mini-0831:phyxminiproblems-problemphyxmini0831-chargeinnanocoulombs}
  \lean{PhyXMiniProblems.ProblemPhyXMini0831.chargeInNanocoulombs}
  \uses{def:physics:phyx-mini-0831:phyxminiproblems-problemphyxmini0831-chargequantity, def:physics:phyx-mini-0831:phyxminiproblems-problemphyxmini0831-chargeincoulombs}
  Read a signed physical charge in nanocoulombs
\end{definition}
\begin{proof}
  This is the declared model definition of charge In Nanocoulombs.
\end{proof}
\begin{definition}[position In Meters]
  \label{def:physics:phyx-mini-0831:phyxminiproblems-problemphyxmini0831-positioninmeters}
  \lean{PhyXMiniProblems.ProblemPhyXMini0831.positionInMeters}
  \uses{def:physics:phyx-mini-0831:phyxminiproblems-problemphyxmini0831-planarpositionquantity}
  Read a planar physical position as Cartesian coordinates in metres
\end{definition}
\begin{proof}
  This is the declared model definition of position In Meters.
\end{proof}
\begin{definition}[field Vector In Newtons Per Coulomb]
  \label{def:physics:phyx-mini-0831:phyxminiproblems-problemphyxmini0831-fieldvectorinnewtonspercoulomb}
  \lean{PhyXMiniProblems.ProblemPhyXMini0831.fieldVectorInNewtonsPerCoulomb}
  \uses{def:physics:phyx-mini-0831:phyxminiproblems-problemphyxmini0831-planarelectricfieldquantity}
  Read a planar electric-field vector in newtons per coulomb
\end{definition}
\begin{proof}
  This is the declared model definition of field Vector In Newtons Per Coulomb.
\end{proof}
\begin{definition}[field Strength In Newtons Per Coulomb]
  \label{def:physics:phyx-mini-0831:phyxminiproblems-problemphyxmini0831-fieldstrengthinnewtonspercoulomb}
  \lean{PhyXMiniProblems.ProblemPhyXMini0831.fieldStrengthInNewtonsPerCoulomb}
  \uses{def:physics:phyx-mini-0831:phyxminiproblems-problemphyxmini0831-electricfieldstrengthquantity}
  Read a nonnegative electric-field strength in newtons per coulomb
\end{definition}
\begin{proof}
  This is the declared model definition of field Strength In Newtons Per Coulomb.
\end{proof}
\begin{definition}[coulomb Constant In Newton Meter Squared Per Coulomb Squared]
  \label{def:physics:phyx-mini-0831:phyxminiproblems-problemphyxmini0831-coulombconstantinnewtonmetersquaredpercoulombsquared}
  \lean{PhyXMiniProblems.ProblemPhyXMini0831.coulombConstantInNewtonMeterSquaredPerCoulombSquared}
  Physlib's Coulomb constant, interpreted as its coherent-SI readout
\end{definition}
\begin{proof}
  This is the declared model definition of coulomb Constant In Newton Meter Squared Per Coulomb Squared.
\end{proof}
\begin{definition}[Charge Source]
  \label{def:physics:phyx-mini-0831:phyxminiproblems-problemphyxmini0831-chargesource}
  \lean{PhyXMiniProblems.ProblemPhyXMini0831.ChargeSource}
  The two point-charge sources shown in the image
\end{definition}
\begin{proof}
  This is the declared model definition of Charge Source.
\end{proof}
\begin{definition}[Figure Point]
  \label{def:physics:phyx-mini-0831:phyxminiproblems-problemphyxmini0831-figurepoint}
  \lean{PhyXMiniProblems.ProblemPhyXMini0831.FigurePoint}
  Distinguished locations shown or determined by the three distance arms
\end{definition}
\begin{proof}
  This is the declared model definition of Figure Point.
\end{proof}
\begin{definition}[Figure Arm]
  \label{def:physics:phyx-mini-0831:phyxminiproblems-problemphyxmini0831-figurearm}
  \lean{PhyXMiniProblems.ProblemPhyXMini0831.FigureArm}
  The three arrows whose labels all read 5.0 cm
\end{definition}
\begin{proof}
  This is the declared model definition of Figure Arm.
\end{proof}
\begin{definition}[Charge Sign]
  \label{def:physics:phyx-mini-0831:phyxminiproblems-problemphyxmini0831-chargesign}
  \lean{PhyXMiniProblems.ProblemPhyXMini0831.ChargeSign}
  Visible charge signs in the two colored circles
\end{definition}
\begin{proof}
  This is the declared model definition of Charge Sign.
\end{proof}
\begin{definition}[figure Point]
  \label{def:physics:phyx-mini-0831:phyxminiproblems-problemphyxmini0831-chargesource-figurepoint}
  \lean{PhyXMiniProblems.ProblemPhyXMini0831.ChargeSource.figurePoint}
  \uses{def:physics:phyx-mini-0831:phyxminiproblems-problemphyxmini0831-chargesource, def:physics:phyx-mini-0831:phyxminiproblems-problemphyxmini0831-figurepoint}
  The location occupied by each point-charge source
\end{definition}
\begin{proof}
  This is the declared model definition of figure Point.
\end{proof}
\begin{definition}[expected Sign]
  \label{def:physics:phyx-mini-0831:phyxminiproblems-problemphyxmini0831-chargesource-expectedsign}
  \lean{PhyXMiniProblems.ProblemPhyXMini0831.ChargeSource.expectedSign}
  \uses{def:physics:phyx-mini-0831:phyxminiproblems-problemphyxmini0831-chargesource, def:physics:phyx-mini-0831:phyxminiproblems-problemphyxmini0831-chargesign}
  The sign visibly associated with each source
\end{definition}
\begin{proof}
  This is the declared model definition of expected Sign.
\end{proof}
\begin{definition}[Opposite Charge Figure]
  \label{def:physics:phyx-mini-0831:phyxminiproblems-problemphyxmini0831-oppositechargefigure}
  \lean{PhyXMiniProblems.ProblemPhyXMini0831.OppositeChargeFigure}
  \uses{def:physics:phyx-mini-0831:phyxminiproblems-problemphyxmini0831-chargesource, def:physics:phyx-mini-0831:phyxminiproblems-problemphyxmini0831-figurearm, def:physics:phyx-mini-0831:phyxminiproblems-problemphyxmini0831-chargesign}
  Literal visual data from image 831.png
\end{definition}
\begin{proof}
  This is the declared model definition of Opposite Charge Figure.
\end{proof}
\begin{definition}[Opposite Point Charge Setup]
  \label{def:physics:phyx-mini-0831:phyxminiproblems-problemphyxmini0831-oppositepointchargesetup}
  \lean{PhyXMiniProblems.ProblemPhyXMini0831.OppositePointChargeSetup}
  \uses{def:physics:phyx-mini-0831:phyxminiproblems-problemphyxmini0831-lengthquantity, def:physics:phyx-mini-0831:phyxminiproblems-problemphyxmini0831-chargequantity, def:physics:phyx-mini-0831:phyxminiproblems-problemphyxmini0831-planarpositionquantity, def:physics:phyx-mini-0831:phyxminiproblems-problemphyxmini0831-planarelectricfieldquantity, def:physics:phyx-mini-0831:phyxminiproblems-problemphyxmini0831-electricfieldstrengthquantity, def:physics:phyx-mini-0831:phyxminiproblems-problemphyxmini0831-chargesource, def:physics:phyx-mini-0831:phyxminiproblems-problemphyxmini0831-figurepoint, def:physics:phyx-mini-0831:phyxminiproblems-problemphyxmini0831-figurearm, def:physics:phyx-mini-0831:phyxminiproblems-problemphyxmini0831-oppositechargefigure}
  The physical sources, their positions, and the fields they produce at the observation dot. No numerical value for the requested strength or preferred answer choice is built into this structure
\end{definition}
\begin{proof}
  This is the declared model definition of Opposite Point Charge Setup.
\end{proof}
\begin{definition}[Matches Primary Figure]
  \label{def:physics:phyx-mini-0831:phyxminiproblems-problemphyxmini0831-matchesprimaryfigure}
  \lean{PhyXMiniProblems.ProblemPhyXMini0831.MatchesPrimaryFigure}
  \uses{def:physics:phyx-mini-0831:phyxminiproblems-problemphyxmini0831-oppositepointchargesetup}
  All unambiguous qualitative information visible in the supplied image
\end{definition}
\begin{proof}
  This is the declared model definition of Matches Primary Figure.
\end{proof}
\begin{definition}[Matches Problem Readouts]
  \label{def:physics:phyx-mini-0831:phyxminiproblems-problemphyxmini0831-matchesproblemreadouts}
  \lean{PhyXMiniProblems.ProblemPhyXMini0831.MatchesProblemReadouts}
  \uses{def:physics:phyx-mini-0831:phyxminiproblems-problemphyxmini0831-lengthincentimeters, def:physics:phyx-mini-0831:phyxminiproblems-problemphyxmini0831-chargeinnanocoulombs, def:physics:phyx-mini-0831:phyxminiproblems-problemphyxmini0831-coulombconstantinnewtonmetersquaredpercoulombsquared, def:physics:phyx-mini-0831:phyxminiproblems-problemphyxmini0831-oppositepointchargesetup}
  The signed charge labels, the three distance labels, and the standard rounded Coulomb constant used by the textbook calculation. None is the requested electric-field strength
\end{definition}
\begin{proof}
  This is the declared model definition of Matches Problem Readouts.
\end{proof}
\begin{definition}[Has Depicted Geometry]
  \label{def:physics:phyx-mini-0831:phyxminiproblems-problemphyxmini0831-hasdepictedgeometry}
  \lean{PhyXMiniProblems.ProblemPhyXMini0831.HasDepictedGeometry}
  \uses{def:physics:phyx-mini-0831:phyxminiproblems-problemphyxmini0831-xaxis, def:physics:phyx-mini-0831:phyxminiproblems-problemphyxmini0831-yaxis, def:physics:phyx-mini-0831:phyxminiproblems-problemphyxmini0831-lengthinmeters, def:physics:phyx-mini-0831:phyxminiproblems-problemphyxmini0831-positioninmeters, def:physics:phyx-mini-0831:phyxminiproblems-problemphyxmini0831-oppositepointchargesetup}
  Cartesian geometry read from the primary figure. The midpoint is chosen as the coordinate origin. The upper and lower charges are one displayed arm above and below it, while the dot is one displayed arm to its left. The metric clauses retain the physical meaning of each arrow independently of the coordinate choice
\end{definition}
\begin{proof}
  This is the declared model definition of Has Depicted Geometry.
\end{proof}
\begin{definition}[Has Physical Parameters]
  \label{def:physics:phyx-mini-0831:phyxminiproblems-problemphyxmini0831-hasphysicalparameters}
  \lean{PhyXMiniProblems.ProblemPhyXMini0831.HasPhysicalParameters}
  \uses{def:physics:phyx-mini-0831:phyxminiproblems-problemphyxmini0831-lengthinmeters, def:physics:phyx-mini-0831:phyxminiproblems-problemphyxmini0831-positioninmeters, def:physics:phyx-mini-0831:phyxminiproblems-problemphyxmini0831-coulombconstantinnewtonmetersquaredpercoulombsquared, def:physics:phyx-mini-0831:phyxminiproblems-problemphyxmini0831-chargesource, def:physics:phyx-mini-0831:phyxminiproblems-problemphyxmini0831-chargesource-figurepoint, def:physics:phyx-mini-0831:phyxminiproblems-problemphyxmini0831-oppositepointchargesetup}
  Positivity and noncoincidence conditions selecting the physical branch
\end{definition}
\begin{proof}
  This is the declared model definition of Has Physical Parameters.
\end{proof}
\begin{definition}[Satisfies Point Charge Electrostatics]
  \label{def:physics:phyx-mini-0831:phyxminiproblems-problemphyxmini0831-satisfiespointchargeelectrostatics}
  \lean{PhyXMiniProblems.ProblemPhyXMini0831.SatisfiesPointChargeElectrostatics}
  \uses{def:physics:phyx-mini-0831:phyxminiproblems-problemphyxmini0831-chargeincoulombs, def:physics:phyx-mini-0831:phyxminiproblems-problemphyxmini0831-positioninmeters, def:physics:phyx-mini-0831:phyxminiproblems-problemphyxmini0831-fieldvectorinnewtonspercoulomb, def:physics:phyx-mini-0831:phyxminiproblems-problemphyxmini0831-fieldstrengthinnewtonspercoulomb, def:physics:phyx-mini-0831:phyxminiproblems-problemphyxmini0831-coulombconstantinnewtonmetersquaredpercoulombsquared, def:physics:phyx-mini-0831:phyxminiproblems-problemphyxmini0831-oppositepointchargesetup}
  Coulomb's vector law in coherent SI units, superposition of the two source fields, and the definition of field strength as the norm of the net field. The law is generic in the two source charges and positions and contains no answer-choice value
\end{definition}
\begin{proof}
  This is the declared model definition of Satisfies Point Charge Electrostatics.
\end{proof}
\begin{theorem}[source to dot distances]
  \label{thm:physics:phyx-mini-0831:phyxminiproblems-problemphyxmini0831-source-to-dot-distances}
  \lean{PhyXMiniProblems.ProblemPhyXMini0831.source_to_dot_distances}
  \uses{def:physics:phyx-mini-0831:phyxminiproblems-problemphyxmini0831-positioninmeters, def:physics:phyx-mini-0831:phyxminiproblems-problemphyxmini0831-oppositepointchargesetup, def:physics:phyx-mini-0831:phyxminiproblems-problemphyxmini0831-matchesproblemreadouts, def:physics:phyx-mini-0831:phyxminiproblems-problemphyxmini0831-hasdepictedgeometry}
  The two charge-to-dot distances implied by the three equal 5.0 cm arms. This is a derived geometric fact, not an assumption of the main answer
\end{theorem}
\begin{proof}
  The conclusion follows from the governing relations and figure readouts named in the statement of source to dot distances.
\end{proof}
\begin{definition}[Answer Choice]
  \label{def:physics:phyx-mini-0831:phyxminiproblems-problemphyxmini0831-answerchoice}
  \lean{PhyXMiniProblems.ProblemPhyXMini0831.AnswerChoice}
  The four answer labels in the problem source
\end{definition}
\begin{proof}
  This is the declared model definition of Answer Choice.
\end{proof}
\begin{definition}[answer Strength In Newtons Per Coulomb]
  \label{def:physics:phyx-mini-0831:phyxminiproblems-problemphyxmini0831-answerstrengthinnewtonspercoulomb}
  \lean{PhyXMiniProblems.ProblemPhyXMini0831.answerStrengthInNewtonsPerCoulomb}
  \uses{def:physics:phyx-mini-0831:phyxminiproblems-problemphyxmini0831-answerchoice}
  The electric-field-strength number printed beside each choice, in N/C
\end{definition}
\begin{proof}
  This is the declared model definition of answer Strength In Newtons Per Coulomb.
\end{proof}
\begin{definition}[recorded Dataset Answer]
  \label{def:physics:phyx-mini-0831:phyxminiproblems-problemphyxmini0831-recordeddatasetanswer}
  \lean{PhyXMiniProblems.ProblemPhyXMini0831.recordedDatasetAnswer}
  \uses{def:physics:phyx-mini-0831:phyxminiproblems-problemphyxmini0831-answerchoice}
  The answer label recorded in the supplied dataset metadata
\end{definition}
\begin{proof}
  This is the declared model definition of recorded Dataset Answer.
\end{proof}
\begin{definition}[answer Agreement Tolerance In Newtons Per Coulomb]
  \label{def:physics:phyx-mini-0831:phyxminiproblems-problemphyxmini0831-answeragreementtoleranceinnewtonspercoulomb}
  \lean{PhyXMiniProblems.ProblemPhyXMini0831.answerAgreementToleranceInNewtonsPerCoulomb}
  Rounding tolerance appropriate to the two-significant-figure answer
\end{definition}
\begin{proof}
  This is the declared model definition of answer Agreement Tolerance In Newtons Per Coulomb.
\end{proof}
\begin{definition}[answer Discrepancy In Newtons Per Coulomb]
  \label{def:physics:phyx-mini-0831:phyxminiproblems-problemphyxmini0831-answerdiscrepancyinnewtonspercoulomb}
  \lean{PhyXMiniProblems.ProblemPhyXMini0831.answerDiscrepancyInNewtonsPerCoulomb}
  \uses{def:physics:phyx-mini-0831:phyxminiproblems-problemphyxmini0831-fieldstrengthinnewtonspercoulomb, def:physics:phyx-mini-0831:phyxminiproblems-problemphyxmini0831-oppositepointchargesetup, def:physics:phyx-mini-0831:phyxminiproblems-problemphyxmini0831-answerchoice, def:physics:phyx-mini-0831:phyxminiproblems-problemphyxmini0831-answerstrengthinnewtonspercoulomb}
  Absolute discrepancy between the calculated strength and a displayed value
\end{definition}
\begin{proof}
  This is the declared model definition of answer Discrepancy In Newtons Per Coulomb.
\end{proof}
\begin{definition}[Is Unique Closest Answer]
  \label{def:physics:phyx-mini-0831:phyxminiproblems-problemphyxmini0831-isuniqueclosestanswer}
  \lean{PhyXMiniProblems.ProblemPhyXMini0831.IsUniqueClosestAnswer}
  \uses{def:physics:phyx-mini-0831:phyxminiproblems-problemphyxmini0831-oppositepointchargesetup, def:physics:phyx-mini-0831:phyxminiproblems-problemphyxmini0831-answerchoice, def:physics:phyx-mini-0831:phyxminiproblems-problemphyxmini0831-answerdiscrepancyinnewtonspercoulomb}
  A displayed choice is strictly closer than every other displayed choice
\end{definition}
\begin{proof}
  This is the declared model definition of Is Unique Closest Answer.
\end{proof}
% --- Archon declaration topology end ---
% --- Archon physics formalization source end ---
